% MINT evaluation paper -- IEEEtran conference draft, claim-centred structure
% (planning/paper_restructure_plan.md, approved 2026-10-27).
% Every result number comes from numbers.tex or tables/*.tex, generated by gen_numbers.py;
% the forest plot comes from figures/forest.py. Do not type results by hand.
% Re-run: python3 paper/gen_numbers.py ; .venv-report/bin/python paper/figures/forest.py
% Wording rule: OPV2V test = primary; Culver City = "unseen town, simulated sensor", secondary; never pooled.
\documentclass[conference]{IEEEtran}
\IEEEoverridecommandlockouts
\usepackage{cite}
\usepackage{amsmath,amssymb}
\usepackage{graphicx}
\usepackage{booktabs}
\usepackage{xcolor}
\usepackage{tikz}
\usetikzlibrary{arrows.meta,positioning,fit,calc}
\usepackage{url}
\usepackage[hidelinks]{hyperref}

\newcommand{\TODO}[1]{\textcolor{red}{[TODO: #1]}}
\newcommand{\CITENEEDED}{\textcolor{red}{[CITATION NEEDED]}}
\newcommand{\pts}{\,pts}
\newcommand{\dev}{\textsuperscript{dev}}
\newcommand{\Cul}{unseen town (simulated sensor)}
% Table bodies are read with the primitive \@@input: since LaTeX 2020-10 \input is robust, and a
% booktabs rule after an \input'ed tabular body fails with "Misplaced \noalign".
\makeatletter
\newcommand{\tabin}[1]{\@@input #1 }
\makeatother

\input{numbers}

\begin{document}

\title{Ray-Verified Tile-Level Counterfactual Gating for Cooperative LiDAR Perception:
A Pre-registered Evaluation in Simulation\\
\thanks{\TODO{funding}}}

\author{\IEEEauthorblockN{\TODO{Author names}}
\IEEEauthorblockA{\TODO{Affiliation}\\
\TODO{City, Country}\\
\TODO{email}}}

\maketitle

\begin{abstract}
Intermediate-fusion cooperative perception lets a malicious sender spoof, remove or shift objects by editing
the features it transmits. We evaluate R1, a defense that asks, per sender and 6.4\,m tile, what changes if the
sender is removed there, and switches the contribution off only when the receiver's own LiDAR rays resolvably
contradict the change. In a pre-registered, one-time test on simulated OPV2V data (\TestScenarios{} scenarios,
\TestFrames{} frames, \NumConditions{} conditions including white-box and B-NORM-aware stealth PGD), the
pre-registered primary, a receiver-only box-level free-space filter (B-RAY), passes its success rule under both
AP definitions, narrowly under global-ranking AP (\VGlobBRAYDelta\pts, one-sided Holm $p=\VGlobBRAYP$, two-sided
CI \VGlobBRAYCI). R1 gains \VGlobRONEDelta\pts{} and is the only defense with a CI above zero in all
\NumConditions{} conditions under global-ranking AP; on the strongest stealth removal$^\dagger$ it restores
\dGlobPgdRemStStrongRONE{} of \LossGlobPgdRemStStrong\pts. Whether R1 beats B-RAY depends on the AP definition
(\VGlobRmB{} vs.\ \VFrameRmB\pts). On a held-out unseen town (Culver City, simulated sensor, \CulFrames{} frames)
the signs of these mean effects repeat, but B-RAY's global-ranking gain is not significant
(\CulHeadGlobBRAY\pts, $p=\CulVGlobBRAYP$). Receiver-blind spoofs stay unrecovered, and R1's compute grows with
the number of cars. Transfer to real sensors is untested.
\end{abstract}

\begin{IEEEkeywords}
cooperative perception, V2V, adversarial robustness, LiDAR, free-space verification, counterfactual,
pre-registration, average precision, simulation
\end{IEEEkeywords}

%======================================================================
\section{Introduction}
%======================================================================
Cooperative perception (CP) lets connected vehicles share what they sense. In \emph{intermediate fusion},
each sender encodes its LiDAR scan into a bird's-eye-view (BEV) feature map, and the receiver fuses the maps
and decodes 3D boxes~\cite{xu2022opv2v,hu2022where2comm,xu2022v2xvit}. The receiver cannot check a feature
map directly, so an insider can perturb what it sends and make the receiver see a car that is not there,
miss one that is, or misplace one~\cite{tu2021adversarial,zhang2024cad,wang2025lucia,zhang2026poseguard}.

Existing defenses either compare outputs with and without a sender and reject whole
agents~\cite{li2023robosac,zhao2024made,hu2025cpguard,tao2026gcp}, test detected boxes against LiDAR
physics~\cite{sun2020carlo,hau2021shadowcatcher,zhang2024cad}, or test feature
statistics~\cite{zhang2026latentguardian}. Agent-level rejection discards a sender's useful information
everywhere to stop an attack that is usually local, and box-level free-space checks can only delete boxes, so
they cannot restore a removed object.

R1 combines the two ideas at a finer grain. It attributes each change in the receiver's output to one sender in
one 6.4\,m tile by decoding with and without that sender, tests the change (an added, removed or shifted box)
against the receiver's own LiDAR rays with a Wilson-95 resolvability margin, and switches the contribution off
only where the receiver can see and the evidence is resolvable. To our knowledge, a tile-level, ray-verified
counterfactual has not been evaluated before; the ray check itself is not new~\cite{sun2020carlo,zhang2024cad}.

We report the evaluation as it was pre-registered, including the parts that do not favour R1. All analyses,
thresholds and attack levels were frozen before the test data were opened, and the test was run once. A rule
fixed on development data made the box-level filter B-RAY the primary and R1 the key secondary
(Appendix~\ref{app:prereg}). After the test, the frozen code was run once on an unseen town of the same
simulator as a secondary check. Every result below is given for the test first and for the unseen town second;
the two are never pooled. All data are simulated.

\textbf{Contributions.}
\begin{enumerate}
\item R1, a tile-level counterfactual defense whose ADD, REMOVE and SHIFT effects are verified against the
receiver's own rays, with UNKNOWN${}={}$keep (Sec.~\ref{sec:method}), and a pre-registered, one-time evaluation
on \TestScenarios{} simulated OPV2V test scenarios with \NumConditions{} conditions, including PGD and
B-NORM-aware stealth PGD, plus a secondary run on an unseen simulated town (Sec.~\ref{sec:protocol}).
\item The finding that the pre-registered primary B-RAY passes narrowly, and that whether R1 improves on it
depends on the AP definition: OPV2V's default frame-order AP reverses the global-ranking statement
(Secs.~\ref{sec:verdict},~\ref{sec:apflip}).
\item The finding that, under global-ranking AP, R1 is the only defense whose CI lies above zero on all four
removal conditions, on both evaluation sets (Sec.~\ref{sec:removal}).
\item Measured limits: receiver-blind spoofs are not recovered, a norm check is evaded by a stealthy attacker,
and R1's clean cost and compute are reported per set (Secs.~\ref{sec:limits},~\ref{sec:costs}).
\end{enumerate}

%======================================================================
\section{Related Work}
%======================================================================
\textbf{Attacks on cooperative perception.}
Tu et al.~\cite{tu2021adversarial} perturb a sender's intermediate message with gradient attacks, the
template for most later work. Zhang et al.~\cite{zhang2024cad} fabricate occlusion-consistent points by ray
casting and run online PGD on features, for spoofing and removal. SOMBRA~\cite{wang2025lucia} hijacks
AttFuse's per-pixel attention on OPV2V to remove objects. PosePert~\cite{zhang2026poseguard} shifts existing
objects by small amounts. The Blind Area Confusion attack~\cite{tao2026gcp} concentrates perturbations where
the victim cannot see, and MVIG~\cite{tao2026mvig} optimizes attack placement against consensus defenses.
TrustFlip~\cite{liu2026trustflip} turns trust-based defenses against benign vehicles. Physical spoofing of the
victim's own LiDAR~\cite{cao2019advlidar} breaks any defense that trusts the receiver's sensor, including ours.

\textbf{Consensus, counterfactual and anomaly defenses.}
ROBOSAC~\cite{li2023robosac} samples sender subsets until the fused result agrees with ego-only perception.
MADE~\cite{zhao2024made} tests each agent with a match loss between outputs with and without that agent and
a reconstruction loss, using conformal p-values and Benjamini--Hochberg control~\cite{benjamini1995fdr}.
CP-Guard~\cite{hu2025cpguard} and GCP~\cite{tao2026gcp} detect malicious agents by consensus, GCP adding
spatial--temporal cues; GLST~\cite{he2026glst} addresses stealthy multi-attacker feature injection.
LUCIA~\cite{wang2025lucia} scales attention by per-agent trust, and MATE~\cite{hallyburton2025mate} maintains
Beta trust over agents and tracks. Latent Guardian~\cite{zhang2026latentguardian} compares collaborator and ego
features by structure, direction and energy. All of these decide per agent (or per subset or track). R1 asks the
same with/without-sender question as MADE and ROBOSAC, but decides per sender and tile, and accepts a decision
only when the receiver's own rays resolve it. We do not compare with these defenses experimentally
(Sec.~\ref{sec:discussion}); our B-NORM baseline is only the simplest member of the feature-statistics family.

\textbf{Free-space defenses.}
CARLO~\cite{sun2020carlo} checks single-vehicle boxes for laser penetration and free space;
Shadow-Catcher~\cite{hau2021shadowcatcher} checks for the LiDAR shadow an object must cast. CAD~\cite{zhang2024cad}
shares occupancy maps between vehicles and runs free-space checks on the fused result, with a local-only
variant. Our B-RAY baseline is a receiver-only box-level check in this style, with a Wilson margin; we claim
no novelty for it.

\textbf{Evaluation practice.}
OPV2V and OpenCOOD~\cite{xu2022opv2v} are the standard simulated benchmark for AttFuse-style intermediate
fusion. Adaptive-attack guidance~\cite{carlini2019evaluating,tramer2020adaptive} asks that defenses be tested
against attackers who know them, which motivates our stealth PGD. AP is usually computed by ranking detections
by score across the whole set, as in PASCAL VOC~\cite{everingham2010voc}; OpenCOOD's default instead accumulates
true and false positives frame by frame. We are not aware of prior CP work that reports both.
Pre-registration~\cite{nosek2018preregistration} is rare in this field.

%======================================================================
\section{Threat Model and Problem Setup}
\label{sec:threat}
%======================================================================
\textbf{System.} A receiver (agent~0) and senders $1,\dots,n{-}1$ ($n\le 5$ on the test set) run AttFuse
PointPillars~\cite{xu2022opv2v,lang2019pointpillars} with the released OPV2V checkpoint, unchanged. Following
OpenCOOD's convention, senders project their points into the receiver's LiDAR frame before encoding, so
sender features $f_{i,k}$ at scale $k\in\{1,2,3\}$ arrive on the receiver's BEV grid. The receiver fuses
them with per-pixel attention and decodes boxes $Y^{+}$.

\textbf{Attacker.} One insider, always sender~1, controls only what it transmits: its features at all three
scales (feature attacks) or its points before encoding (point attack). For PGD the attacker is white-box:
it knows the weights, the receiver's features, the fusion and the heads. It does not control the receiver's
sensor, pose or software, and it does not collude with others. No attack is optimized against R1's or B-RAY's
ray test; the receiver-blind spoof (Sec.~\ref{sec:attacks}) is the only attack aware of where the receiver can
see. Frames with $n=1$ have no sender and are scored as clean.

\textbf{Defender.} The receiver trusts only its own LiDAR scan $P_0$ and pose. It may decode received features
any number of times. It never uses ground truth at inference.

\textbf{Goal.} Maximize AP@0.7 for cars on attacked frames while losing as little as possible on clean
frames. A defense may delete boxes (B-RAY), drop senders (B-NORM, ego-only) or switch off a sender in a tile
(R1).

%======================================================================
\section{Method}
\label{sec:method}
%======================================================================
\input{figures/r1_method.tex}

\subsection{Counterfactual effects}
Let $Y^{+}$ be the receiver's boxes decoded from all agents and $Y^{-}(i)$ the boxes decoded with sender
$i$ excluded (Fig.~\ref{fig:r1method}, steps 1--2):
\begin{equation}
\begin{split}
Y^{+}&=\mathrm{Dec}\big(\mathrm{Fuse}(f_0,f_1,\dots,f_{n-1})\big),\\
Y^{-}(i)&=\mathrm{Dec}\big(\mathrm{Fuse}(\{f_j\}_{j\neq i})\big).
\end{split}
\label{eq:cf}
\end{equation}
We match $Y^{+}$ and $Y^{-}(i)$ by Hungarian assignment on BEV IoU (match if IoU${}\ge 0.1$) and define the
effects of sender $i$:
\begin{equation}
\begin{split}
\mathrm{ADD}_i&=\{b\in Y^{+} : b \text{ unmatched in } Y^{-}(i)\},\\
\mathrm{REMOVE}_i&=\{b\in Y^{-}(i) : b \text{ unmatched in } Y^{+}\},\\
\mathrm{SHIFT}_i&=\{(b^{+},b^{-})\in M_i : d_c>0.3\,\text{m} \lor d_\psi>10^\circ\},
\end{split}
\label{eq:effects}
\end{equation}
where $M_i$ is the set of matched pairs, $d_c=\|c^{+}-c^{-}\|$ the distance of the BEV box centres and $d_\psi=|\psi^{+}-\psi^{-}|$ the heading
difference. Each effect belongs to the 6.4\,m tile $u$ that contains its box centre (in either version for
SHIFT). The \emph{unit} of decision is the pair $(i,u)$.

\subsection{Receiver ray evidence}
For a box $b$, each return of the receiver's own scan $P_0$ defines a ray from the receiver's sensor. A ray that
ends inside $b$ is \emph{support} (SUP); one that passes through $b$ and ends behind it is \emph{traversal}
(TRAV); one that ends before $b$ carries no information. Rays without a return are never counted as
traversal. The support ratio is
\begin{equation}
s(b)=\frac{\mathrm{SUP}(b)}{N(b)},\qquad N(b)=\mathrm{SUP}(b)+\mathrm{TRAV}(b).
\label{eq:s}
\end{equation}
An effect is KNOWN if every box it involves has $N\ge 3$ and is not the receiver's own box; otherwise its unit
is UNKNOWN.

\subsection{Resolvability gate}
For $k$ supporting rays out of $N$, the Wilson score interval~\cite{wilson1927} with $z=1.96$ is
\begin{equation}
[\mathrm{lo},\mathrm{hi}]=\frac{k+\tfrac{z^2}{2}\mp z\sqrt{\tfrac{k(N-k)}{N}+\tfrac{z^2}{4}}}{N+z^2}.
\label{eq:wilson}
\end{equation}
Rays through one box are spatially correlated, so this is a \emph{resolvability margin} on counts, not a
probability. Each KNOWN effect receives a verdict, with $s^{+}$ ($s^{-}$) the ratio of its box in $Y^{+}$
($Y^{-}(i)$), $I^{\pm}$ their Wilson intervals, and $\mathrm{dis}$ the event $I^{+}\cap I^{-}=\emptyset$:
\begin{equation}
\begin{array}{@{}lll@{}}
 & \text{gate iff} & \text{useful iff}\\[1pt]
\text{ADD:} & \mathrm{hi}(s^{+})<0.5 & \mathrm{lo}(s^{+})>0.5\\
\text{REMOVE:} & \mathrm{lo}(s^{-})>0.5 & \mathrm{hi}(s^{-})<0.5\\
\text{SHIFT:} & \mathrm{dis}\land s^{+}<s^{-} & \mathrm{dis}\land s^{+}>s^{-}
\end{array}
\label{eq:gate}
\end{equation}
The ``useful'' conditions mirror the gate conditions: an effect the rays resolvably confirm protects its unit.

\subsection{Tile gating and composed output}
With $E_{i,u}$ the effects of sender $i$ in tile $u$, the unit is gated iff it is KNOWN, at least one of its
effects gates, and none is resolvably useful:
\begin{equation}
\begin{split}
g_{i,u}=\mathbf{1}\big[&\mathrm{KNOWN}_{i,u}\land \exists e{\in}E_{i,u}{:}\ \mathrm{gate}(e)\\
&\land \nexists e{\in}E_{i,u}{:}\ \mathrm{useful}(e)\big].
\end{split}
\label{eq:unitgate}
\end{equation}
UNKNOWN units are kept. Let $S_u=\{i: g_{i,u}=1\}$, and let $C_u(S)$ be the pre-NMS candidates in tile $u$ of the
decode without the senders in $S$, so that $C_u(\emptyset)$ are those of $Y^{+}$. The output is composed tile by
tile:
\begin{equation}
Y_{\mathrm{R1}}=\mathrm{NMS}\Big(\textstyle\bigcup_{u} C_u(S_u)\Big).
\label{eq:compose}
\end{equation}
This re-uses the decodes of~\eqref{eq:cf}; a new decode is needed only for tiles that gate two or more senders.
Composition was checked against a direct masked decode on every test and unseen-town scenario (all identical).
R1 has no tunable threshold. The ablation \emph{R1-noREMOVE} never lets a REMOVE verdict gate (its mirror
condition is kept).

\subsection{Baselines}
\textbf{B-RAY} (pre-registered primary) applies R1's ADD test to the final boxes without any counterfactual:
\begin{equation}
\begin{split}
Y_{\mathrm{B\text{-}RAY}}&=\{b\in Y^{+} : \lnot\,\kappa(b)\},\\
\kappa(b)&=\mathrm{KNOWN}(b)\land \mathrm{hi}\big(s(b)\big)<0.5.
\end{split}
\label{eq:bray}
\end{equation}
It is a receiver-only free-space filter in the style of CARLO's penetration check and CAD's local-only
variant~\cite{sun2020carlo,zhang2024cad}. It can only delete boxes. Because of the Wilson margin, a box needs at
least four traversing rays and no support to be deleted.

\textbf{B-NORM} drops sender $i$ for the whole frame if, at any scale $k$,
$\overline{|f_{i,k}|}/\overline{|f_{0,k}|}>\tau_k$, with $\tau_k$ the 99.5th percentile of this ratio over clean
development frames (block~H; $\tau_1=\BnormTauOne$, $\tau_3=\BnormTauThree$). \textbf{Ego-only} drops all
senders. No baseline has a parameter fitted at test time.

\subsection{Ground-truth harm labels (evaluation only)}
To audit what R1 switches off, each unit is labelled offline with ground truth. With Hungarian matching at
BEV IoU${}\ge 0.5$ over the boxes centred in tile $u$,
\begin{equation}
\begin{split}
L_u(Y)&=\#\mathrm{FP}+\#\mathrm{FN}+\tfrac12\textstyle\sum_{\mathrm{TP}}(1-\mathrm{IoU}),\\
D_{i,u}&=L_u(Y^{+})-L_u\big(Y^{-}(i,u)\big).
\end{split}
\label{eq:harm}
\end{equation}
A unit is HARMFUL if $D_{i,u}>0.5$, USEFUL if $D_{i,u}<-0.5$ and NEUTRAL otherwise. Gating a USEFUL unit is a
\emph{false switch-off}. The labels are never used by any defense.

%======================================================================
\section{Evaluation Protocol}
\label{sec:protocol}
%======================================================================
\subsection{Two evaluation sets}
\label{sec:sets}
We use OPV2V~\cite{xu2022opv2v}, simulated in CARLA~\cite{dosovitskiy2017carla}, with the released AttFuse
checkpoint on CPU (Table~\ref{tab:splits}). \textbf{The OPV2V test} is the primary, one-time evaluation: the
OPV2V test split minus one scenario opened earlier, \TestScenarios{} scenarios and \TestFrames{} frames ($n=1$
in \TestNOne{} frames, $n=2$ in \TestNTwo, $n=3$ in \TestNThree{} and $n\ge 4$ in \TestNFourPlus). Before the
test only clean detector outputs on it had been seen. \textbf{Culver City} is the secondary evaluation: OPV2V's
separate town, simulated with the same CARLA LiDAR model, with a layout modelled on that city, and unseen during
detector training and during all of our development (\CulScenarios{} scenarios, \CulFrames{} frames,
$n=\CulNmin$--\CulNmax). We call it the \emph{\Cul}. It was planned and committed before its data were
downloaded (D-72) and run once, after the test, with methods, $\sigma_k$, $\tau_k$, attack levels, seeds and
analysis unchanged; nothing was fitted on it. A planned check on its first scenario, before any defense or attack
ran, gave clean AP@0.7 of \CulSanityAP{} (frame order, \CulSanityFrames{} frames); over all its scenarios the
no-defense frame-order AP@0.7 is \CulAPCleanFrame{} (\CulPublishedAP{} released with the
checkpoint~\cite{xu2022opv2v}). The two sets are reported side by side and never pooled; with four scenarios,
the unseen town's scenario CIs are coarse. The development blocks T and H (Table~\ref{tab:splits}) appear only in
the history of Appendices~\ref{app:prereg} and~\ref{app:devtest}.

\begin{table}[!t]
\caption{Data blocks. The OPV2V test is the primary, one-time evaluation; Culver City (unseen town, simulated
sensor) is a secondary evaluation after the test. T and H are development blocks.}
\label{tab:splits}
\centering
\footnotesize
\setlength{\tabcolsep}{4pt}
\begin{tabular}{@{}lrrcp{4.05cm}@{}}
\toprule
\tabin{tables/tab_splits.tex}
\bottomrule
\end{tabular}
\end{table}

\subsection{Attack conditions}
\label{sec:attacks}
Table~\ref{tab:attacks} lists the \NumConditions{} conditions. The copy-paste families (spoof, remove, shift)
edit 3.2\,m-aligned patches of the attacker's feature maps at random locations; \emph{spoof (points)} copies a
car's points to an empty spot and re-encodes; \emph{spoof, receiver-blind} places the spoof only where the
receiver has fewer than three rays. An attack is \emph{applied} when it can be placed and \emph{realized} when it
changes the undefended output; unapplied frames are scored as clean.

\textbf{PGD.} The attacker perturbs its features at all scales within $\|\delta_k\|_\infty\le\epsilon_k=r\,\sigma_k$,
with $\sigma_k$ the standard deviation of clean sender features on block~H. With step $\eta_k=\epsilon_k/4$, 10
sign-gradient ascent steps~\cite{madry2018pgd} and a uniform random start,
\begin{equation}
\delta_k^{t+1}=\Pi_{\mathcal{B}_k}\!\big(\delta_k^{t}+\eta_k\,\mathrm{sign}\,\nabla_{\delta_k}\mathcal{L}(\delta^{t})\big),
\label{eq:pgd}
\end{equation}
where $\mathcal{B}_k=\{\delta:|\delta|\le\epsilon_k,\ f_{1,k}+\delta\ge 0\}$. The loss uses the clean fused output as
pseudo-labels (no ground truth): \emph{untargeted} maximizes classification and regression loss on all anchors;
\emph{remove} minimizes the scores of the clean boxes. The \textbf{stealth} variants add B-NORM's constraint after
every step through an exact projection that keeps every scale's magnitude ratio below $0.95\,\tau_k$
(Appendix~\ref{app:stealth}); frames where no such perturbation exists are not attacked and are scored as clean.
The stealth families were designed on block~H after the plain PGD results there (D-68), and their levels
$r=1.0$ (marked $^\dagger$) were added after seeing development results (D-70); the $^\dagger$ conditions are
never used for a criterion.

\begin{table}[!t]
\caption{Attack conditions: frames on which the attack changed the undefended output (Real., of
\WithSenderSpoofFeat{} test and \CulWithSender{} Culver City frames with a sender) and attack loss in AP@0.7
points under global-ranking (G) and frame-order (F) AP. PGD levels are $r$; st.\ = B-NORM-aware stealth PGD.
$^\dagger$Added after seeing development results (D-70).}
\label{tab:attacks}
\centering
\footnotesize
\setlength{\tabcolsep}{2.6pt}
\begin{tabular}{@{}lrrrrrr@{}}
\toprule
\tabin{tables/tab_attacks_both.tex}
\bottomrule
\end{tabular}
\end{table}

\subsection{Metrics}
$\Delta$AP@0.7 is the AP@0.7 of a defense minus that of no defense on the same frames, in AP points. With the
attack loss $L=\mathrm{AP}(\text{clean, none})-\mathrm{AP}(\text{attacked, none})$,
\begin{equation}
\begin{split}
\mathrm{Rec}&=\Delta\mathrm{AP}/L\qquad(\text{n/a if }L<0.005),\\
\mathrm{CL}_m&=1-\mathrm{AP}^{\mathrm{clean}}_{m}\,/\,\mathrm{AP}^{\mathrm{clean}}_{\mathrm{none}},
\end{split}
\label{eq:rec}
\end{equation}
where $\mathrm{CL}_m$ is the relative clean loss of defense $m$ and $\mathrm{AP}^{\mathrm{clean}}_{m}$ its AP@0.7
on clean frames.
\textbf{Two AP definitions.} Let detections $j$ carry a score $c_j$, a frame index $t_j$ and a true-positive
flag $\mathrm{tp}_j$ (IoU${}\ge 0.7$), and let $G$ be the number of ground-truth cars. For an ordering $\pi$,
\begin{equation}
\begin{split}
p_k&=\tfrac{1}{k}\textstyle\sum_{j\le k}\mathrm{tp}_{\pi(j)},\qquad
r_k=\tfrac{1}{G}\textstyle\sum_{j\le k}\mathrm{tp}_{\pi(j)},\\
\mathrm{AP}_\pi&=\textstyle\sum_k (r_k-r_{k-1})\,\max_{k'\ge k}p_{k'}.
\end{split}
\label{eq:ap}
\end{equation}
\emph{Global-ranking AP} (G, our headline) uses $\pi_G$, all detections sorted by descending $c_j$.
\emph{Frame-order AP} (F), the OpenCOOD default used in OPV2V papers, uses $\pi_F$, which concatenates detections
frame by frame in order of $t_j$. Under $\pi_F$, a false positive in an early frame lowers the precision of every
later detection, whatever its score. Our NumPy implementation matches OpenCOOD's to $\NumpyOpencoodDiff$.

\subsection{Statistics and labels}
CIs come from a scenario bootstrap~\cite{efron1979bootstrap} ($B=\BootB$ resamples, seed 0; 15 clusters on the
test, 4 on the unseen town); a 10-frame block bootstrap stratified by scenario is secondary. For a statistic
$\theta$ with bootstrap replicates $\theta^*_b$,
\begin{equation}
p=\frac{1+\#\{b:\theta^*_b\le 0\}}{B+1}.
\label{eq:p}
\end{equation}
B-RAY and R1 are co-tested against no defense with Holm's step-down procedure~\cite{holm1979} at $\alpha=0.05$
within each AP definition: with $p_{(1)}\le p_{(2)}$, reject $H_{(1)}$ if $p_{(1)}\le\alpha/2$, then $H_{(2)}$ if
also $p_{(2)}\le\alpha$. The pre-registered success rule requires B-RAY to be Holm-rejected on the mean
$\Delta$AP@0.7 over the four original attacks (spoof, remove, shift by feature; spoof by points) \emph{and} the
upper scenario-CI bound of its clean loss to be at most 1\%, under \emph{both} AP definitions (D-64). The key
secondary, R1~$-$~B-RAY, is a paired difference whose scenario CI selects one of three pre-written statements.

We label every result. \textbf{Confirmatory}: the success rule, the Holm co-test and the R1~$-$~B-RAY statement,
on the OPV2V test only. \textbf{Descriptive}: everything else, including every unseen-town number (there the
success rule is applied unchanged, but it does not re-open the test verdict). \textbf{Post hoc} ($^\dagger$):
choices made after seeing development results. Development numbers appear only in Appendix~\ref{app:devtest}.

%======================================================================
\section{Results}
\label{sec:results}
%======================================================================
Each subsection states one claim and then gives the OPV2V test (primary) and the unseen town (secondary).

\subsection{The pre-registered verdict}
\label{sec:verdict}
\textbf{The primary B-RAY passes on the test, narrowly under global-ranking AP} (Table~\ref{tab:verdict},
Fig.~\ref{fig:forest}). On the OPV2V test, B-RAY improves the mean of the four original attacks by
\VGlobBRAYDelta\pts{} under global-ranking AP: the one-sided bootstrap test rejects after Holm ($p=\VGlobBRAYP$,
threshold \VGlobBRAYHolmThr), while the two-sided 95\% scenario CI \VGlobBRAYCI{} includes zero. Under frame-order
AP the gain is \VFrameBRAYDelta\pts{} \VFrameBRAYCI{} ($p=\VFrameBRAYP$). Its clean loss stays inside the 1\%
bound under both definitions, so the success rule passes under both (confirmatory). On the \Cul, the same rule,
applied descriptively, does not pass under global-ranking AP (\CulHeadGlobBRAY\pts{} \CulHeadGlobBRAYCI,
$p=\CulVGlobBRAYP$) and passes under frame-order AP (\CulHeadFrameBRAY\pts{} \CulHeadFrameBRAYCI). The unseen
town's status is therefore ``AP-definition-dependent''; this does not change the test verdict.

\textbf{R1, co-tested,} gains \VGlobRONEDelta\pts{} \VGlobRONECI{} under global-ranking AP (rejected,
$p=\VGlobRONEP$) and \VFrameRONEDelta\pts{} \VFrameRONECI{} under frame-order AP (not rejected, $p=\VFrameRONEP$)
on the test. On the unseen town it is rejected under both (\CulHeadGlobRONE{} \CulHeadGlobRONECI{} and
\CulHeadFrameRONE\pts{} \CulHeadFrameRONECI). R1 has no success rule of its own.

The effects are small in absolute terms because the four original attacks cost little over all frames
(test: \LossGlobSpoofFeat, \LossGlobRemoveFeat, \LossGlobShiftFeat{} and \LossGlobSpoofPts\pts{} under global-ranking
AP; Table~\ref{tab:attacks}).

\begin{table*}[!t]
\caption{Pre-registered verdict on the OPV2V test (confirmatory) and the same rule applied descriptively on
Culver City (unseen town, simulated sensor; 4 scenarios, coarse CIs). Mean $\Delta$AP@0.7 over the four original
attacks, AP points, 95\% scenario-bootstrap CI, one-sided bootstrap $p$. Never pooled. $^\ast$Pre-registered
statement: the counterfactual ``improves over'' or ``performs worse than'' the box-level filter, or its benefit
``is not established''.}
\label{tab:verdict}
\centering
\footnotesize
\setlength{\tabcolsep}{6pt}
\begin{tabular}{@{}lcccc@{}}
\toprule
\tabin{tables/tab_verdict_both.tex}
\bottomrule
\end{tabular}
\end{table*}

\begin{figure}[!t]
\centering
\includegraphics[width=\columnwidth]{figures/forest.pdf}
\caption{Mean $\Delta$AP@0.7 over the four original attacks (B-RAY, R1, R1-noREMOVE) and the paired
R1~$-$~B-RAY, OPV2V test (primary, \TestScenarios{} scenarios, filled) and Culver City (unseen town, simulated
sensor; secondary, \CulScenarios{} scenarios, open), under both AP definitions. Bars: 95\% scenario-bootstrap CI.
Labels: Holm decision (co-tested methods) or the pre-registered R1~$-$~B-RAY statement. Never pooled; B-NORM and
ego-only are off scale (Appendix~\ref{app:full}).}
\label{fig:forest}
\end{figure}

\subsection{Removal attacks}
\label{sec:removal}
\textbf{Under global-ranking AP, R1 is the only defense whose scenario CI lies above zero on all four removal
conditions, on both sets} (Table~\ref{tab:conditions}; descriptive). On the OPV2V test, R1 gains
\dGlobRemoveFeatRONE\pts{} \ciGlobRemoveFeatRONE{} on copy-paste removal, \dGlobPgdRemRONE\pts{}
\ciGlobPgdRemRONE{} on PGD removal and \dGlobPgdRemStRONE\pts{} \ciGlobPgdRemStRONE{} on stealth PGD removal; on
the \Cul{} the same three gains are \CulDGlobRemoveFeatRONE{} \CulCiGlobRemoveFeatRONE,
\CulDGlobPgdRemRONE{} \CulCiGlobPgdRemRONE{} and \CulDGlobPgdRemStRONE\pts{} \CulCiGlobPgdRemStRONE. B-RAY's
effect on these conditions is zero or negative on both sets: it can only delete boxes.

On the strongest stealth removal$^\dagger$, the attack costs \LossGlobPgdRemStStrong\pts{} on the test and R1
restores \dGlobPgdRemStStrongRONE{} \ciGlobPgdRemStStrongRONE{} (recovery \RecGlobPgdRemStStrongRONE; frame-order
\dFramePgdRemStStrongRONE{} \ciFramePgdRemStStrongRONE, recovery \RecFramePgdRemStStrongRONE). On the unseen town
it costs \CulLossGlobPgdRemStStrong\pts{} and R1 restores \CulDGlobPgdRemStStrongRONE{}
\CulCiGlobPgdRemStStrongRONE{} (recovery \CulRecGlobPgdRemStStrongRONE; frame-order
\CulDFramePgdRemStStrongRONE, recovery \CulRecFramePgdRemStStrongRONE). At that level ego-only also gains
(\dGlobPgdRemStStrongEgo{} and \CulDGlobPgdRemStStrongEgo\pts), because the attack costs more than cooperation is
worth.

Under frame-order AP, R1's removal gain is resolved only at the strong level$^\dagger$ on the test (on
copy-paste removal it is \dFrameRemoveFeatRONE\pts{} \ciFrameRemoveFeatRONE), and on the two PGD removals at $r=0.25$
(\CulDFramePgdRemRONE{} \CulCiFramePgdRemRONE{} and \CulDFramePgdRemStRONE\pts{} \CulCiFramePgdRemStRONE) and the
strong level on the unseen town. Without REMOVE verdicts (R1-noREMOVE) the strong-level gain is
\dGlobPgdRemStStrongRNoRem\pts{} on the test and \CulDGlobPgdRemStStrongRNoRem\pts{} on the unseen town, so the
REMOVE test carries the recovery.

\begin{table*}[!t]
\caption{$\Delta$AP@0.7 vs.\ no defense per condition for B-RAY and R1, OPV2V test (primary) and Culver City
(unseen town, simulated sensor; secondary), both AP definitions, in AP points. \textbf{Bold}: 95\%
scenario-bootstrap CI above zero; \textit{italic}: CI below zero. CIs and the other defenses:
Tables~\ref{tab:full-g} and~\ref{tab:full-f}. $^\dagger$Added after seeing development results (D-70).
Never pooled.}
\label{tab:conditions}
\centering
\footnotesize
\setlength{\tabcolsep}{4pt}
\begin{tabular}{lcccccccc}
\toprule
\tabin{tables/tab_conditions.tex}
\bottomrule
\end{tabular}
\end{table*}

\subsection{Every condition}
\label{sec:every}
\textbf{Under global-ranking AP, R1 is the only defense whose scenario CI lies above zero in all
\NumConditions{} conditions of the test, clean frames included} (Table~\ref{tab:conditions}; two of the conditions
are $^\dagger$). On the \Cul{} its point estimate is positive in all \NumConditions, and its CI lies above zero in
all but two: clean frames (\CulDGlobCleanRONE\pts{} \CulCiGlobCleanRONE) and the receiver-blind spoof
(\CulDGlobSpoofBlindRONE\pts{} \CulCiGlobSpoofBlindRONE). Under frame-order AP this does not hold on either set:
on the test R1 loses on clean frames (\dFrameCleanRONE\pts{} \ciFrameCleanRONE) and its CIs include zero on most
attacks, and on the unseen town its clean estimate is \CulDFrameCleanRONE\pts. B-RAY's gains are concentrated on
spoofing under frame-order AP (test \dFrameSpoofFeatBRAY, unseen town \CulDFrameSpoofFeatBRAY\pts{} on feature
spoofs).

\subsection{R1 versus B-RAY, and the AP-definition flip}
\label{sec:apflip}
\textbf{Whether R1 improves on B-RAY depends on the AP definition} (Fig.~\ref{fig:forest}, bottom rows). On the
OPV2V test, R1~$-$~B-RAY is \VGlobRmB\pts{} \VGlobRmBCI{} under global-ranking AP, for which the pre-registered
statement reads ``\VGlobRmBStatement'', and \VFrameRmB\pts{} \VFrameRmBCI{} under frame-order AP:
``\VFrameRmBStatement'' (confirmatory). On the \Cul, the global-ranking statement is the same (\CulVGlobRmB\pts{}
\CulVGlobRmBCI); the frame-order difference has the same sign but is not resolved (\CulVFrameRmB\pts{}
\CulVFrameRmBCI: ``\CulVFrameRmBStatement''). The same sign pattern was seen on development data and decided the
choice of the primary (Appendix~\ref{app:prereg}).

The two definitions differ in where a low-score false positive is charged: at the tail of the global ranking,
where it costs little, or at its frame's position in the frame-order accumulation, where it lowers the precision
of every later frame. Consistent with this, the same feature spoof costs \LossFrameSpoofFeat{} vs.\
\LossGlobSpoofFeat\pts{} on the test. Deleting low-score fakes, which B-RAY does, and R1's switch-offs on clean
frames (Sec.~\ref{sec:costs}) are charged under the same rules. \textbf{A ranking between CP defenses can thus
depend on an evaluation setting that OPV2V papers rarely state}, and we report both.

\subsection{What is not defended}
\label{sec:limits}
\textbf{The receiver-blind spoof is not recovered.} No defense recovers more than 0.10 of its loss on either set
under either AP definition, and at most 0.06 on the frames where the spoof changed the output. R1's all-frame
recovery is \RecGlobSpoofBlindRONE{} / \RecFrameSpoofBlindRONE{} on the test and \CulRecGlobSpoofBlindRONE{} /
\CulRecFrameSpoofBlindRONE{} on the unseen town (global / frame-order); realized-only it is
\RealRecGlobSpoofBlindRONE{} / \RealRecFrameSpoofBlindRONE{} and \CulRealRecGlobSpoofBlindRONE{} /
\CulRealRecFrameSpoofBlindRONE. B-RAY's is at most \RecFrameSpoofBlindBRAY{} (test) and
\CulRecFrameSpoofBlindBRAY{} (unseen town). On the test, \BlindRealizedLowRay{} of the \RealizedSpoofBlind{}
realized spoofs landed where the receiver had fewer than three rays (or on its own box) and
\BlindRealizedKnownHit{} touched a KNOWN unit; on the unseen town \CulBlindRealizedLowRay{} of
\CulRealizedSpoofBlind{} and \CulBlindRealizedKnownHit. R1's global-ranking gain on this condition on the test
(\dGlobSpoofBlindRONE\pts) is of the same size as its gain on clean frames (\dGlobCleanRONE\pts). The
visibility-limited attacker oracle bounds recovery at \VisOracleRecSpoofBlind{} on the test (frame-order; the
unconstrained oracle reaches \OracleRecSpoofBlind).

\textbf{B-NORM is evaded by the stealth attacker.} B-NORM flags the attacker on \CatchPgdUnt{} of plain PGD frames
on the test and \CulCatchPgdUnt{} on the unseen town, and on \CatchPgdUntSt{} of frames under every
B-NORM-aware stealth condition on both sets, while the stealth removal at $r=0.25$ still costs
\LossGlobPgdRemSt\pts{} (test) and \CulLossGlobPgdRemSt\pts{} (unseen town). On clean frames B-NORM flags the
benign sender on \CatchClean{} (test) and \CulCatchClean{} (unseen town) of sender frames, against
\BnormFalseFlagH{} on the development block where its threshold was fitted; it changes the mean of the four
original attacks by \HeadGlobBNORM{} and \CulHeadGlobBNORM\pts{} (global-ranking AP).

\textbf{B-RAY recovers little of optimized attacks.} B-RAY recovers at most \RecFramePgdUntStStrongBRAY{} (test)
and \CulRecFramePgdUntBRAY{} (unseen town) of any PGD attack's loss and gains nothing on PGD removal on either set
(Table~\ref{tab:conditions}); its frame-order gains on untargeted PGD have CIs above zero on both sets. R1's
recovery of untargeted PGD is also small (at most \RecGlobPgdUntStRONE{} under global-ranking AP on the test).

\subsection{Costs}
\label{sec:costs}
\textbf{Clean cost} (Table~\ref{tab:clean}). B-RAY's clean relative loss stays inside the pre-registered 1\%
bound under both AP definitions on the test (\CleanGlobBRAY, \CleanFrameBRAY; confirmatory as part of the rule)
and on the \Cul{} (\CulCleanGlobBRAY, \CulCleanFrameBRAY). On clean test frames it drops \BrayCleanDrop{} of
\BrayCleanBoxes{} boxes, \BrayCleanDropGT{} of them matched to cars. R1's clean cost depends on the AP definition:
\CleanGlobRONE{} under global ranking and \CleanFrameRONE{} \CleanFrameRONECI{} under frame order on the test, where
it switches off \RoneCleanGated{} of \RoneCleanUnits{} units, \textbf{\RoneCleanFalseSwitch{} of them useful}
(\RoneCleanFalseSwitchRemove{} via REMOVE; R1-noREMOVE: \RNoRemCleanFalseSwitch). On the unseen town the
frame-order cost is \CulCleanFrameRONE{} \CulCleanFrameRONECI, with \CulRoneCleanGated{} units gated and
\CulRoneCleanFalseSwitch{} false switch-offs. B-NORM costs \CleanGlobBNORM{} and \CulCleanGlobBNORM{} under
global-ranking AP.

\textbf{Compute} (Table~\ref{tab:runtime}). Dense R1's overhead over receiver-only perception grows with the
number of cars because it decodes once per sender: on the test \OvDenseTwoPfifty, \OvDenseThreePfifty{} and
\OvDenseFourPlusPfifty{} (p50) at $n=2$, 3 and ${\ge}4$; on the unseen town \CulOvDenseTwoPfifty{} and
\CulOvDenseThreePfifty{} at $n=2$ and 3 (\TimingMachine, CPU only). The cheaper trigger T(0.05) is a cost
variant only: it was not exact on development data, so no accuracy number uses it. The project's original 10\%
p95 budget is not met by any R1 variant; we report cost and do not judge it (a post hoc change, D-50).

\begin{table}[!t]
\caption{Clean cost: relative AP@0.7 loss on clean frames (95\% scenario CI below each value) under
global-ranking (G) and frame-order (F) AP. Lower block, what each defense removes on clean frames: B-RAY, boxes
dropped / boxes (matched to a ground-truth car); R1 and R1-noREM, units gated / units (USEFUL, Eq.~\eqref{eq:harm});
B-NORM, benign sender flagged / sender frames.}
\label{tab:clean}
\centering
\scriptsize
\setlength{\tabcolsep}{3.5pt}
\begin{tabular}{@{}lcccc@{}}
\toprule
\tabin{tables/tab_clean_both.tex}
\bottomrule
\end{tabular}
\end{table}

\begin{table}[!t]
\caption{Overhead per frame vs.\ receiver-only perception, p50 / p95 in \%, by number of agents $n$ (\TimingMachine,
CPU only). Base: receiver perception p50 in seconds. Dense and T(0.05): R1; Culver City has $n\le 3$.}
\label{tab:runtime}
\centering
\footnotesize
\setlength{\tabcolsep}{4pt}
\begin{tabular}{@{}llrcccc@{}}
\toprule
\tabin{tables/tab_runtime_both.tex}
\bottomrule
\end{tabular}
\end{table}

%======================================================================
\section{Discussion and Limitations}
\label{sec:discussion}
%======================================================================
\textbf{Receiver-blind regions.} Every receiver-ray defense keeps what the receiver cannot see, and R1 likewise
cannot verify a removal there. Closing this gap needs evidence other than the receiver's own rays, such as other
benign senders' scans, which brings back the trust problem.

\textbf{Adaptive attackers.} Stealth PGD shows how cheaply a defense that only looks at feature magnitudes is
bypassed. No attack was optimized against R1's or B-RAY's ray test: an attacker who places fakes where the rays are
ambiguous (fewer than four traversing rays) or that the receiver cannot see is not stopped. The attack levels,
the primary and the stealth attacks were chosen on exposed development data (block~H), and the strong levels after
seeing results there. The test was run once; its results cannot tune anything further without new data.

\textbf{Baselines.} We did not run learned or consensus defenses (e.g.\ LUCIA, PoseGuard, MADE); B-NORM stands in
only for the simplest feature-statistics check, not for those methods, and we make no claim relative to them.

\textbf{Single attacker, single detector.} The attacker is always sender~1, alone. Two colluding senders could
each confirm the other's fake; R1's per-sender counterfactual removes one sender at a time and would miss it. R1's
exact composition relies on AttFuse's per-pixel fusion; detectors that mix information spatially after fusion
(e.g.\ V2X-ViT~\cite{xu2022v2xvit}) would need the tile locality re-verified.

\textbf{Simulation.} OPV2V is CARLA simulation~\cite{dosovitskiy2017carla}. The unseen town tests transfer to a
new town layout under the same simulated sensor, not to real sensors, real traffic or real attacks. Real LiDAR
misses dark or glass surfaces (more UNKNOWN) and returns from vegetation or spray (false support). Real-sensor
datasets such as V2V4Real~\cite{xu2023v2v4real} and Mixed Signals~\cite{luo2025mixedsignals} were not evaluated.
The receiver's sensor is assumed honest; physical spoofing~\cite{cao2019advlidar} would break that.

\textbf{Compute.} Dense R1 scales with the number of senders and exceeds the receiver's own perception time at
$n\ge 4$ on the test and at $n=3$ on the unseen town on our CPU.

\textbf{Statistics.} The bootstrap $p$ of Eq.~\eqref{eq:p} is the bootstrap mass at or below zero, not a
null-calibrated test, and percentile intervals over 15 (test) or 4 (unseen town) unequal clusters may
under-cover. Attacks that are not realized and $n=1$ frames dilute every all-frame $\Delta$AP equally. The
calibrated-probability estimand of our original plan was dropped before the test (Appendix~\ref{app:prereg}).

%======================================================================
\section{Conclusion}
%======================================================================
In a pre-registered, one-time test on simulated data, a receiver-only box-level free-space filter (B-RAY) passed
its pre-registered rule, narrowly under global-ranking AP; on an unseen simulated town its global-ranking gain was
not significant. Whether R1, a tile-level counterfactual verified by the receiver's own rays, improves on B-RAY
depends on the AP definition, so CP papers should report both. Under global-ranking AP, R1 was the only defense
whose CI lay above zero on all four removal conditions, on both sets. Receiver-blind spoofs were not recovered by
any defense, and a norm check was evaded by a stealthy attacker. Future work: a hybrid that uses B-RAY's box filter
for added objects and only R1's REMOVE test for removals (untested), attacks optimized against the ray test, and
validation on real-sensor data never touched by method design.

\section*{Data and Code Availability}
\TODO{repository URL, frozen commit, pre-registration files and checksums.}

%======================================================================
% Appendices (IEEE order: before the acknowledgment and the references)
%======================================================================
\appendices

\section{Pre-registration history}
\label{app:prereg}
Every rule was committed before the data it governs were opened. Block~T was used to write R1; block~H was a
held-out check of the frozen R1 (E4) and was exposed afterwards. The original pre-registration (v1) made R1 the
primary. Before any B-RAY result on block~H existed, rule D-57 fixed: if the block-H 95\% CI of R1~$-$~B-RAY
(frame-order AP, 10-frame bootstrap) had its lower bound above zero, R1 would stay primary; otherwise B-RAY would
become primary and R1 the key secondary. On H the difference was negative (\DevHFrameRmB\pts{} \dev), so B-RAY
became the primary (D-58). Pre-registration v2 added the Holm co-test of B-RAY and R1 against no defense (D-59).
A pre-test audit then found that the sign of R1~$-$~B-RAY on H depends on the AP definition (D-60;
Table~\ref{tab:devtest}). Rather than switch the primary a second time, v3 (D-64) requires the v2 success rule to
pass under \emph{both} AP definitions, quotes headline numbers in global-ranking AP, and reports
``AP-definition-dependent'' where the two disagree. v3 also added B-NORM and the PGD conditions as secondary
conditions (D-65--D-68); the stealth families were designed on H after the plain PGD results there (D-68). v3.1
(D-70) added the two strong stealth levels ($r=1.0$, $^\dagger$) \textbf{after seeing development results on H};
they are never used for a criterion. v3.2 (D-71) fixed a bug in the stream code that crashed the first attempt
before any result was written; no method, condition or criterion changed. The original 10\% p95 latency budget
was dropped as a pass/fail claim post hoc (D-50), and the calibrated-probability estimand of the charter was
dropped when its calibration block became block~H (D-52); the Wilson interval is a resolvability margin, not a
probability (D-40). After the test, the unseen-town evaluation was planned and committed before any of its data
were downloaded (D-72) and recorded as a secondary, descriptive result (D-73).

\section{Development versus test}
\label{app:devtest}
Table~\ref{tab:devtest} places the test next to the development blocks. These are \textbf{development numbers,
not test results}. The blocks were smaller (T: \DevTScenarios{} scenarios, H: \DevHScenarios), H had only
$n=\DevHN$, and B-RAY was selected on H. Gains are smaller on the test than on H under both definitions; the sign
pattern of R1~$-$~B-RAY is the same on H and on the test.

\begin{table}[!t]
\caption{Development (T in-sample, H exposed) vs.\ OPV2V test: mean $\Delta$AP@0.7 over the four original
attacks in AP points with 95\% scenario CI; G = global-ranking AP, F = frame-order AP. Development numbers are not
test results.}
\label{tab:devtest}
\centering
\scriptsize
\setlength{\tabcolsep}{2.2pt}
\begin{tabular}{@{}lccc@{}}
\toprule
\tabin{tables/tab_dev_vs_test.tex}
\bottomrule
\end{tabular}
\end{table}

\section{Full per-condition results}
\label{app:full}
Tables~\ref{tab:full-g} and~\ref{tab:full-f} give every defense in every condition with its 95\%
scenario-bootstrap CI, for each set and AP definition.

\begin{table*}[!t]
\caption{Every defense in every condition, \textbf{global-ranking AP}: $\Delta$AP@0.7 vs.\ no defense, AP points,
95\% scenario-bootstrap CI. Bold: CI above zero; italic: CI below zero. Culver City is the unseen town (simulated
sensor), 4 scenarios, coarse CIs; the two sets are never pooled. PGD levels are $r$; st.\ = B-NORM-aware stealth
PGD. $^\dagger$Added after seeing development results (D-70).}
\label{tab:full-g}
\centering
\scriptsize
\setlength{\tabcolsep}{3pt}
\begin{tabular}{@{}lccccc@{}}
\toprule
\tabin{tables/tab_full_global.tex}
\bottomrule
\end{tabular}
\end{table*}

\begin{table*}[!t]
\caption{As Table~\ref{tab:full-g}, \textbf{frame-order AP} (OpenCOOD default).}
\label{tab:full-f}
\centering
\scriptsize
\setlength{\tabcolsep}{3pt}
\begin{tabular}{@{}lccccc@{}}
\toprule
\tabin{tables/tab_full_frame.tex}
\bottomrule
\end{tabular}
\end{table*}

\section{Stealth projection and expected outcomes}
\label{app:stealth}
After every PGD step (and after the random start), each scale $k$ is projected onto the box
$[\ell_k,\epsilon_k]$ with $\ell_k=\max(-\epsilon_k,-f_{1,k})$, intersected with B-NORM's mean-magnitude bound.
With $m_k(\lambda)$ the mean of $f_{1,k}+\mathrm{clip}(\delta_k-\lambda,\ell_k,\epsilon_k)$, the exact Euclidean
projection is
\begin{equation}
\begin{split}
\delta_k&\leftarrow \mathrm{clip}\big(\delta_k-\lambda_k,\ \ell_k,\ \epsilon_k\big),\\
\lambda_k&=\min\big\{\lambda\ge 0:\ m_k(\lambda)\le 0.95\,\tau_k\,\overline{|f_{0,k}|}\big\},
\end{split}
\label{eq:stealth}
\end{equation}
with $\lambda_k$ found by bisection. Frames where the constraint has no solution are not attacked (``stealth
infeasible'', scored as clean): on the test \InfeasiblePgdUntSt, \InfeasiblePgdRemSt, \InfeasiblePgdUntStStrong{}
and \InfeasiblePgdRemStStrong{} of \WithSenderPgdRemSt{} sender frames for the four stealth conditions; on the
unseen town \CulInfeasiblePgdUntSt, \CulInfeasiblePgdRemSt, \CulInfeasiblePgdUntStStrong{} and
\CulInfeasiblePgdRemStStrong{} of \CulWithSender.

Expectations written before any PGD run were checked descriptively. ``B-RAY recovery $<0.10$ under both AP
definitions'' held in five of six PGD conditions on the test (it failed on strong untargeted stealth
PGD$^\dagger$, frame-order recovery \RecFramePgdUntStStrongBRAY) and in four of six on the unseen town (it failed
on plain untargeted PGD, \CulRecFramePgdUntBRAY, and on strong untargeted stealth PGD$^\dagger$,
\CulRecFramePgdUntStStrongBRAY). ``R1's CI above zero under both definitions'' on removal held on the test only at
the strong level$^\dagger$, and on the unseen town on all three PGD removal conditions. ``B-NORM catch
${\le}5\%$'' held for all stealth conditions on both sets.

\section{Agent count}
\label{app:agents}
Under frame-order AP (10-frame CIs), B-RAY's mean gain on the test falls from \ByNTwoBRAY\pts{} \ByNTwoBRAYCI{}
at $n=2$ to \ByNFourPlusBRAY\pts{} \ByNFourPlusBRAYCI{} at $n\ge 4$, and R1's from \ByNTwoRONE{} to
\ByNFourPlusRONE\pts{} ($n\ge 4$ covers \ByNFourPlusScen{} scenarios). On the unseen town B-RAY gains
\CulByNTwoBRAY\pts{} \CulByNTwoBRAYCI{} at $n=2$ (\CulByNTwoScen{} scenarios) and \CulByNThreeBRAY\pts{}
\CulByNThreeBRAYCI{} at $n=3$ (\CulByNThreeScen{} scenario), and R1 \CulByNTwoRONE{} and \CulByNThreeRONE\pts.
These strata are descriptive.

\section*{Acknowledgment}
\TODO{acknowledgments}

\bibliographystyle{IEEEtran}
\bibliography{refs}

\end{document}
